\chapter{Technical Background}
\label{ch:bg}

This chapter fixes the vocabulary used in the rest of the thesis. Nothing
here is a contribution; everything here is required in order to read the
implementation chapters without a pile of footnotes.

\section{Controller Area Network}
\label{sec:can}

The Controller Area Network is a multi-master, message-oriented serial bus
originally specified by Bosch and now standardised as
ISO~11898~\cite{bosch1991can,iso11898,voss2008can}. Nodes do not address
each other; they publish \emph{frames} that carry an identifier, a data
payload and a cyclic redundancy check. Arbitration is bitwise and
non-destructive: the identifier with the most dominant bits wins, so
priority is encoded in the identifier itself. That property is the reason
automotive and industrial control systems still prefer CAN over
point-to-point links when a modest amount of data must be shared among many
nodes with bounded latency~\cite{navet2005can}.

A classical CAN 2.0 frame carries up to 8~bytes of payload and an 11-bit
(standard) or 29-bit (extended) identifier. CAN FD~\cite{iso11898fd} keeps
the same arbitration phase and then switches to a higher bit-rate for the
data phase, with payloads of up to 64~bytes. The SR5E1 Stellar-E
microcontroller implements an FDCAN peripheral that is source-compatible
with both. In this project the physical layer was operated as CAN~2.0 with
11-bit identifiers: the diagnostic and command payloads used here fit in
eight bytes, and the dyno software at the far end of the cable was
configured for that mode. The firmware drivers, however, go through the
FDCAN HAL, so a later migration to CAN FD is a configuration change rather
than a protocol redesign.

Two protocol facts matter for everything that follows.

\begin{enumerate}[leftmargin=1.6em]
  \item \textbf{The identifier \emph{is} the message.} There is no
        connection handshake and no session. A node that wants to start the
        motor publishes identifier \texttt{0x33}; a node that wants to
        report heatsink temperature publishes identifier \texttt{0x5A}.
        Receivers filter on identifiers. This is why a DBC file, which
        binds names and layouts to identifiers, is not documentation on the
        side---it is the only human-readable contract the bus has.
  \item \textbf{Payload layout is convention, not syntax.} CAN does not
        know about integers, physical units or byte order. Little-endian
        packing of a 32-bit speed in RPM, a 16-bit current in
        \si{\milli\ampere}, or a 32-bit PI word that concatenates a gain
        and a divisor, is entirely a matter of agreement between producer
        and consumer. Break that agreement and the bus still ``works'':
        frames are acknowledged at the data-link layer, while the
        application reads garbage.
\end{enumerate}

\section{CAN databases (DBC)}
\label{sec:dbc-bg}

A DBC file is the de-facto industry dictionary for a CAN
bus~\cite{vectordbc}. For every message it records

\begin{itemize}[leftmargin=1.6em]
  \item the identifier and the data length code (DLC);
  \item a list of signals, each with start bit, length, byte order,
        signedness, factor, offset, range and unit;
  \item optional value tables, multiplexors and comments.
\end{itemize}

Tools such as Vector CANalyzer~\cite{vectorcanalyzer} compile the DBC and
use it in both directions: incoming frames are decoded into physical
signals (``heatsink is \SI{47.3}{\celsius}''), and outgoing signals are
encoded into frames. If the DBC is wrong, the GUI is wrong, even when the
firmware is right. Chapter~\ref{ch:dbc} is, in that sense, as important as
the C changes in Chapter~\ref{ch:dev}.

A recurring design choice in this project is to treat the firmware as
authoritative and the DBC as a generated-looking projection of it. Offsets
of $+$\texttt{0x30} and $+$\texttt{0x40}, described in
\cref{ch:arch}, exist precisely so that a reader of either artefact can
predict the other.

\section{Field-oriented control of PMSM machines}
\label{sec:foc}

Permanent-magnet synchronous machines are controlled, in the ST Motor
Control SDK as in most industrial drives, by field-oriented control
(FOC)~\cite{park1929,leonhard1996foc,krishnan2010pmsm}. Stator currents
are transformed into a rotor-aligned $dq$ frame in which

\begin{itemize}[leftmargin=1.6em]
  \item $I_d$ is the flux-producing current,
  \item $I_q$ is the torque-producing current.
\end{itemize}

In the linear region, electromagnetic torque is proportional to $I_q$
(and, with reluctance torque, to a cross term in $I_d I_q$). Flux weakening
at high speed is obtained by commanding a negative $I_d$. The two currents
are regulated by inner PI loops; an outer PI loop on mechanical speed (or
on torque) produces the $I_q$ reference. PWM of the inverter legs
then synthesises the corresponding stator voltages~\cite{holmes2003pwm}.

Three consequences for the CAN interface follow directly from this
structure.

\begin{enumerate}[leftmargin=1.6em]
  \item Independent $I_d$ and $I_q$ references are not a luxury. Characterisation
        of the machine---identification of flux maps, of $L_d$ and $L_q$, of
        iron-loss maps---requires holding one axis while sweeping the other.
        A protocol that can only write the pair $(I_d,I_q)$ atomically is
        already useful; a protocol that can also ramp each axis on its own
        schedule is what a dyno actually needs.
  \item Ramps are a safety and a physics feature. A step of $I_q$ is a step
        of torque; on a dyno with a finite mechanical bandwidth it produces
        a torsional shock. Speed ramps and current ramps, parameterised by
        a final value and a duration in milliseconds, are therefore
        first-class commands, not cosmetic.
  \item PI gains are dimensioned as a numerator and a power-of-two
        divisor~\cite{astromhagglund2006}. The MCSDK stores
        $K_p$ as an integer together with
        \texttt{KPDivisorPOW2}, so that the effective gain is
        $K_p / 2^{n}$. Exposing only $K_p$ over CAN would make the
        closed-loop behaviour impossible to reconstruct; exposing both in
        two messages would make atomic updates impossible. The packing
        described in \cref{sec:pi-pack} is the resulting compromise.
\end{enumerate}

\section{The Stellar-E SR5E1 platform}
\label{sec:stellar}

Stellar-E is STMicroelectronics' automotive microcontroller family built
around Arm Cortex-R52 cores, with a functional-safety architecture aimed at
ISO~26262 applications~\cite{ststellar}. The SR5E1 device used on the bench
provides one or more FDCAN controllers, a rich timer set for PWM
generation, analog peripherals for current and voltage sensing, and a debug
port that Lauterbach TRACE32 can attach to. From the point of view of this
thesis the relevant properties are pedestrian and essential:

\begin{itemize}[leftmargin=1.6em]
  \item incoming CAN frames raise an interrupt
        (\texttt{IRQ\_CAN1\_LINE0\_HANDLER} in the project sources);
  \item outgoing frames are submitted to the FDCAN TX FIFO from a software
        handler (\texttt{DBCCAN\_TX\_IRQ\_Handler} /
        \texttt{CFCP\_TX\_IRQ\_Handler});
  \item the Motor Control SDK is linked as a set of weakly-defined
        \texttt{\_\_weak} functions, which is what made it possible to add
        \texttt{UI\_ExecFluxRamp} and the PI packing without forking the
        whole library.
\end{itemize}

The ``dishwasher'' qualifier that appears throughout the code is historical.
ST's Motor Control SDK ships application profiles---among them an appliance
profile originally demonstrated on a dishwasher drum motor. The register
map, the command codes and the Frame Communication Protocol (FCP) used by
that profile were reused as the internal language of the dyno interface.
DYNO2DW is the adapter between that language and the identifier map that
the bench already spoke.

\section{ST Motor Control protocol and FCP}
\label{sec:mcp}

The internal protocol is a short, length-prefixed, checksummed frame:

\begin{center}
\texttt{[Code : 1 byte][Length : 1 byte][Payload : Length bytes][CRC : 1 byte]}
\end{center}

\texttt{Code} selects an operation (\texttt{EXECUTE\_CMD}, \texttt{SET\_REG},
\texttt{GET\_REG}, set-ramp, set-current, \ldots). \texttt{Length} is the
payload size, not the total frame size. The CRC is computed by
\texttt{SelfCRCCalc} over the preceding bytes. On the MCU the reconstructed
bytes are pushed, one by one, into
\texttt{CFCP\_RX\_IRQ\_Handler}, which implements the Frame Communication
Protocol state machine, including inter-byte timeouts. On the way out,
\texttt{CFCP\_TX\_IRQ\_Handler} serialises a response: either a 1-byte
acknowledgement or a 4-byte register value.

Two observations, both of which will reappear as bugs in
\cref{ch:val,ch:dev}, follow from this layout.

\begin{itemize}[leftmargin=1.6em]
  \item The CAN payload is \emph{not} identical to the FCP frame. DYNO2DW
        is a translator: it reads a CAN identifier and a raw 8-byte buffer,
        and it \emph{builds} the FCP frame in a local array \texttt{msg[]}.
        Any confusion between ``byte~0 of the CAN payload'' and ``byte~0 of
        the FCP frame'' produces a protocol discrepancy. That is exactly
        the SET\_REG case study of \cref{sec:discrepancy}.
  \item Acknowledgements and register-read responses are both produced by
        the same TX handler. If they share a CAN identifier, a consumer
        cannot tell a successful write from a four-byte value. Separating
        them by identifier, using frame size as the discriminant, is the
        repair described in \cref{sec:ack}.
\end{itemize}

\section{Dynamometer test benches}
\label{sec:dyno}

A dynamometer imposes a mechanical boundary condition on a rotating
machine: a speed, a torque, or a more elaborate profile such as a driving
cycle~\cite{isermann2014engine}. In an inverter characterisation campaign
the dyno software is typically the master of the experiment. It

\begin{itemize}[leftmargin=1.6em]
  \item brings the UUT to a defined state (idle, aligned, running);
  \item programs references (speed, $I_q$, $I_d$) and the time over which
        they must be reached;
  \item samples electrical and thermal quantities in lock-step with the
        mechanical ones;
  \item reacts to faults (over-current, over-voltage, over-temperature)
        faster than a human operator.
\end{itemize}

All four items are CAN messages in this project. The first is
\texttt{EXECUTE\_COMMAND}. The second is the family of set-point and ramp
commands. The third is the ten-message diagnostic stack emitted every
\SI{100}{\milli\second}. The fourth is the combination of the machine-state
diagnostic, the fault-acknowledge flag and the NACK path.

The GUI developed in \cref{ch:gui} does not replace the dyno software. It
replaces the \emph{manual} use of PCAN-View during bring-up and during
tests that a human still wants to drive---PI tuning, visual inspection of
current loops, fault recovery---while using the same DBC that a fully
automated CAPL script, or the dyno computer itself, can consume.

\section{Tooling}
\label{sec:tools}

Three tools dominate the validation chapters and are therefore introduced
here, briefly.

\paragraph{PCAN-View} is PEAK-System's bus monitor and injector, used with
a PCAN-USB FD adapter~\cite{pcanview}. It is the lowest-level stimulus:
the engineer types an identifier, a DLC and eight hex bytes, and the
adapter puts that frame on the wire. It is also the lowest-level oracle:
if the ten diagnostic identifiers do not appear every \SI{100}{\milli\second},
the physical layer or the TX path is broken, regardless of what the GUI
says.

\paragraph{TRACE32} is Lauterbach's hardware debugger~\cite{trace32}.
Attached to the SR5E1 debug port it can halt the core on the FDCAN
interrupt, step through \texttt{DYNO2DWProtocol}, and show the local
\texttt{msg[]} buffer and the Motor Control handles. Combined with PCAN-View
it gives an end-to-end trace from ``frame on the bus'' to ``function
called in the SDK''.

\paragraph{Vector CANalyzer} sits one level above both~\cite{vectorcanalyzer}.
It loads the DBC, runs CAPL~\cite{vectorcapl} on identifier and system-variable
events, and exposes those system variables to any COM client. The Python GUI
of \cref{ch:gui} is one such client. CANalyzer, not Python, owns the
real-time path; Python owns widgets, gauges and plots.
